% !PRINTED 153
\weberpage{153}
%<<CONT
durch die $x$ ausgedrückt mittelst einer neuen linearen Substitution $E$, die wir so bezeichnen können:

\weq{10}{z = E (x) = B A (x),}

und die Substitution $E = BA$ heisst aus $B$ und $A$ \gesperrt{zusammengesetzt} oder \gesperrt{componirt}. Es ist dabei aber zwischen $AB$ und $BA$ zu unterscheiden. Zwei Substitutionen $A, B$ von der besonderen Eigenschaft, dass $AB = BA$ ist, heissen mit einander \gesperrt{vertauschbar} oder \gesperrt{commutativ}. Die Substitutionscoëfficienten von $E$ ergeben sich durch Einsetzen der Ausdrücke (1) in (9):

% [sic: Weber prints i below the first sum although h is its summation index; 1,n is expected.]
\weq{11}{\begin{gathered} z_k = \sum_{i}^{h} e_h^{(k)} x_h, \\ e_h^{(k)} = \sum_{1,n}^{i} b_i^{(k)} a_h^{(i)}. \end{gathered}}

Diese Formeln sind ganz dieselben, die wir im §. 27 des ersten Bandes benutzt haben\footnote{Nur war dort aus einem leicht ersichtlichen Grunde die Bezeichnung etwas anders.}, und wir können demnach das in (11) ausgedrückte Gesetz der Composition der Substitutionen so ausdrücken:

\gesperrt{1. Um die aus zwei Substitutionen $B, A$ zusammengesetzte Substitution $BA$ zu bilden, verfährt man ganz so, als ob die beiden Determinanten $|B|, |A|$ nach der Multiplicationsregel mit einander multiplicirt werden sollten. Es sind dabei, um die Elemente einer Zeile zu bilden, die Elemente einer Zeile der ersten Componente mit den entsprechenden Elementen der Colonnen der zweiten Componente zu multipliciren und dann zu addiren. Man drückt dies auch kurz so aus, dass in der ersten Componente nach Zeilen, in der zweiten nach Colonnen summirt wird.}

Aus dieser Regel ergiebt sich die Folgerung:

\gesperrt{2. Die Determinante einer zusammengesetzten Substitution ist gleich dem Producte aus den Determinanten der Componenten.}
